% TOP_PROBLEM: {"id": 2639, "title": "Kourovka Notebook Problem 21.130", "category_id": 17, "category": {"id": 17, "name": "group_theory", "display_name": "Group Theory", "description": "Problems about groups, group actions, representations, and related algebraic structures.", "slug": "group-theory", "order_index": 17, "created_at": "2026-07-31T15:26:25.671Z"}, "status": "open", "problem_number": "KOU-21.130", "set_id": 10}
% FIRST_POSTED: 2026-10-01
% ENTRY_KIND: statement_only
% UnsolvedMath ID 2639; source catalogue code KOU-21.130
% !TeX program = lualatex
% Definition and sources only; this file is not an LLM solution attempt.
\documentclass[11pt,a4paper]{article}
\usepackage[margin=25mm]{geometry}
\usepackage{fontspec}
\setmainfont{lmroman10-regular.otf}[BoldFont=lmroman10-bold.otf,
 ItalicFont=lmroman10-italic.otf,BoldItalicFont=lmroman10-bolditalic.otf]
\usepackage{amsmath,amssymb,mathtools}
\usepackage{unicode-math}
\setmathfont{latinmodern-math.otf}
\AtBeginDocument{\let\setminus\smallsetminus}
\usepackage{xurl,hyperref}
\hypersetup{colorlinks=true,urlcolor=blue}
\setcounter{secnumdepth}{0}
\setlength{\parindent}{0pt}
\setlength{\parskip}{5pt}
\setlength{\emergencystretch}{3em}
\begin{document}
\section*{Kourovka Notebook Problem 21.130}\label{2639}
{\small UnsolvedMath ID 2639 · KOU-21.130 · Group Theory · Kourovka Notebook - New Problems, Issue 21}
\subsection{Definitions and mathematical statement}
Let $G$ be a finite abelian group, written additively with identity $0$, and suppose $|G|$ is odd. Let $A\subseteq G$ be a subset of cardinality $n\ge3$. An ordering of $A$ is a list $(a_1,\ldots,a_n)$ containing each member exactly once. Interpret subscripts cyclically by setting $a_{n+1}=a_1$.

The conjecture asks whether there is always an ordering for which
\[
 a_i+a_{i+1}\ne a_j+a_{j+1}
 \qquad(1\le i<j\le n).
\]
Thus all $n$ adjacent sums around the cycle must be distinct, including the closing sum $a_n+a_1$. The sums are computed in $G$, so they use its group operation and any torsion relations of $G$.

Equivalently, consider the complete graph on the elements of $A$, assigning to an edge $\{a,b\}$ the label $a+b\in G$. The required ordering is a Hamilton cycle whose edge labels are all different. This restatement uses undirected edges because $G$ is abelian.

The subset is arbitrary and need not generate $G$, contain $0$, or be a subgroup. Its adjacent sums need not belong to $A$ and may include $0$; only repetition among them is forbidden. The conclusion is existential for each individual subset, and cyclic rotations or reversal of a successful ordering are allowed. Oddness refers to the ambient group order, not to the size $n$ of the subset.
\subsection{Short English statement}
Can every subset of size at least three in a finite abelian group of odd order be cyclically ordered with all neighboring sums distinct?
\subsection{Sources}
\begin{itemize}
\item[S1] ULAM.AI and the UnsolvedMath Contributors, supplied problems.json, record 2639 (KOU-21.130), Kourovka Notebook - New Problems, Issue 21. Catalogue identity and source transcription; CC BY 4.0.
\url{https://huggingface.co/datasets/ulamai/UnsolvedMath}.
\item[S2] The Kourovka Notebook, 21st edition, version 47 (30 September 2026), new problems, Problem 21.130. Expanded formulation of the supplied catalogue statement.
\url{https://arxiv.org/pdf/1401.0300v47}.
\end{itemize}
\end{document}
